\section{Attribution in detail}\label[appendix]{sec:attribution}

Throughout the paper four levels of ownership are distinguished.
\emph{Classical} tools are used as published. \emph{Imported}
statements are restated, credited to their origins and proved in
\cref{app:earlier}. \emph{Adapted} statements are restated because
a hypothesis changes, and the change is proved here. \emph{New}
components are the constructions and applications introduced in this
paper, rather than the standard identities used to prove them.
\Cref{tab:attribution} assigns each
component of the proof to one of these levels.

Joint Lean companion: \leanrepo.

The companion's source mapping documents how each printed algorithm of~\cite{BCGT13} is transcribed.

\paragraph{The algorithmic method and the transfer.}
The algorithm-to-lower-bound program is Williams'~\cite{Wil13}. The
almost-everywhere transfer used here follows
Chen--Lyu--Williams~\cite{CLW20}, including the constructive
refuter, which adapts the almost-everywhere hierarchy of
Fortnow and Santhanam~\cite{FS16}, and the general-circuit hardness
amplification of Sudan, Trevisan, and Vadhan~\cite{STV01}, in the form
of~\cite[Lemma~3.9]{CLW20}. The transfer also uses the PCP
refinements of Chen and Williams~\cite{CW19} and of Bathie and
Williams~\cite{BW24},
the proximity construction of Ben-Sasson, Goldreich, Harsha, Sudan, and
Vadhan~\cite{BGHSV06}, and the short PCP in the form stated by Murray
and Williams~\cite[Section~5]{MW20}, credited there to Ben-Sasson and
Viola~\cite{BV14}. The author's preprint~\cite{NagCG} assembles
these into a transfer from certified pair counting to
almost-everywhere constant-error hardness, through a weighted
analysis of the MAX-2SAT gadget of
Garey--Johnson--Stockmeyer~\cite{GJS76}; its Appendix~D proves the
short PCP in the form it uses, from the construction of Ben-Sasson,
Chiesa, Genkin, and Tromer~\cite{BCGT13}. The full transfer is given here with the wire measure; its
supporting constructions are reproduced in \cref{app:earlier}. We claim no new PCP, amplification or
hierarchy theorem.

\paragraph{Relation to~\cite{Nag72}.}
The author's preprint~\cite{Nag72} introduced the partition into
products of Hamming-code stars~\cite{Hamming}, the mixed-residual
population bound, and the universal fixed-moment sampler with exact
short-seed conditional expectations, whose concentration step is
Tropp's matrix Bernstein inequality~\cite[Theorem~1.6]{Tropp}. Its
counter applies the sampler to Boolean state matrices built from
pointwise count lists and \(\F_2\) filters, and it pays roughly the
square root of the population. We keep its cells and its sampler and
replace everything in between: the sampled matrices are Fourier
characters of the residual, the potential is computed by walk-sum
tomography instead of count lists, and the masks are exact integer
polynomials instead of \(\F_2\) filters. No count list, radix state,
polarization or Boolean lift of~\cite{Nag72} is used. We also import
its separated exact printer and the lemmas that feed it.

\paragraph{Relation to~\cite{NagCG}.}
The author's preprint~\cite{NagCG} reconstructs exact score
distributions of restricted circuits by nested subsampling, moment
certificates and modular tomography, building on Hajdu and
Tijdeman's factorization method for discrete tomography~\cite{HT01}.
We apply the same machinery to a different integer, the signed
residual sum of a closed walk; the localization lemma is re-proved
for walk sums, whose summands are bounded by \(2m\) rather than one.
We also use its transfer, as described above.

\paragraph{Classical tools.}
Integer normalization of real threshold functions is classical,
due to Muroga, Toda, and Takasu~\cite{MTT61}, in the form
of~\cite[Lemma~3.2]{CTW26}. The decomposition of a threshold function
into disjoint exact thresholds goes back to Hansen and
Podolskii~\cite{HP10}; the first-crossing construction used is that
of Chen and Williams~\cite[Proposition~18(2) and Appendix~B]{CW19}.
Batch evaluation uses Williams' rectangular matrix
multiplication~\cite[Corollary~4.4]{WilACC14}, based on Coppersmith's
algorithm~\cite{Cop82}, with Matou\v{s}ek's bucketing~\cite{Mat91}.
The modulus-amplifying polynomial is that of Beigel and
Tarui~\cite{BT94}. Exact rational linear systems are solved by
fraction-free elimination~\cite{Bareiss68}. Newton's forward-difference
formula, the Chernoff and Bernstein bounds, and finite Fourier
analysis are standard.



\begin{table}[!htbp]
\centering\footnotesize
\begin{tabularx}{\textwidth}{@{}>{\raggedright\arraybackslash}p{.30\textwidth}>{\raggedright\arraybackslash}p{.10\textwidth}>{\raggedright\arraybackslash}X@{}}
\toprule
Component (where used) & Level & Source and change\\
\midrule
Signed-top gate corollary (\cref{thm:gates}) & Classical & Replication and finite-cube threshold complementation; wire cost at most \((n+1)\topwt(C)\)\\
Product-star partition (\cref{imp:partition}) & Imported & \cite[Lemma~A.1]{Nag72}; Hamming codes~\cite{Hamming}\\
Range of the mixed residual (\cref{imp:range}) & Imported & \cite[Lemma~A.2]{Nag72}; we use only the integer form of \(\Delta\), never its \(\F_2\) form\\
Mixed-residual population (\cref{imp:population}, \cref{lem:meanpop}) & Imported & \cite[Theorem~A.4]{Nag72}; the constant \(6\) is computed here\\
Score identity, mixed dependence (\cref{lem:identity,lem:mixed}) & Adapted & \cite[Lemma~A.2 and the proof of Theorem~A.4]{Nag72}, restated for two circuits with an explicit pointwise bound\\
Fourier decomposition with selectors (\cref{lem:fourier}) & Adapted & Standard finite Fourier expansion; \cite[Equation~(29)]{Nag72}; here also for parity calls\\
Trace bound for \([-1,1]\) entries (\cref{lem:tracebound}) & Adapted & The computation of~\cite[Lemma~F.1]{Nag72} for \(0/1\) entries, redone; concentration by~\cite[Theorem~1.6]{Tropp}\\
Uniform selector error, potential to error (\cref{lem:selector,lem:root}) & Adapted & \cite[Lemma~F.1 and its proof]{Nag72}, for matrices with entries in \([-1,1]\) and a subset of columns\\
Short seed and exact conditional expectations (\cref{imp:seed}) & Imported & \cite[Lemma~F.2 and its proof]{Nag72}\\
Cycle-sum identity, conditional potentials (\cref{lem:cycle,cor:conditional}) & New & \\
Exact integer mask (\cref{def:mask,lem:maskpoly}) & New & Newton forward differences; replaces the \(\F_2\) filters of~\cite[Appendix~D]{Nag72}\\
Mask guarantees (\cref{lem:smaskpop,lem:filter,lem:smaskmass}) & New & Chernoff bounds\\
Nested rates and telescoping (\cref{imp:chain}) & Imported & \cite[Lemma~5.1 and Equation~(1)]{NagCG}\\
Walk-sum localization (\cref{lem:localize}) & Adapted & \cite[Lemma~5.2]{NagCG} for walk sums: summands bounded by \(2m\)\\
Moment certificates (\cref{lem:moments}) & Imported & \cite[Lemma~5.3]{NagCG}, any index set\\
Certified boxes (\cref{lem:certify}) & Adapted & \cite[Sections~7.3 and~8]{NagCG}, for walk sums and the exact mask\\
One- and two-score reconstruction (\cref{imp:window}) & Imported & \cite[Lemmas~4.1--4.4]{NagCG}; Hajdu--Tijdeman~\cite{HT01}\\
Bernstein inequality (\cref{imp:bernstein}) & Classical & \cite[Lemma~3.4]{NagCG}, from~\cite[Theorem~2.8.4]{Vershynin18}\\
Fermat tests and lifting (\cref{imp:amp}, \cref{lem:residue}) & Imported / adapted & \cite[Lemmas~6.6 and~6.7]{NagCG}; Beigel--Tarui~\cite{BT94}; here with signed walk weights and the exact mask\\
Integer normalization (\cref{imp:integer}) & Classical & \cite[Lemma~3.2]{NagCG} \(=\) \cite[Lemma~3.2]{CTW26}; \cite{MTT61}; applied on each gate's own inputs\\
Disjoint exact children, pullback, stacking (\cref{imp:disjoint,imp:pullback,imp:stack}) & Imported & \cite[Lemmas~C.2, C.3, C.6]{Nag72}; \cite{HP10,CW19}\\
Separated printer (\cref{imp:printer}) & Imported & \cite[Theorem~C.5]{Nag72}; Williams~\cite{WilACC14}, Coppersmith~\cite{Cop82}, Matou\v{s}ek~\cite{Mat91}\\
Printing walk polynomials (\cref{lem:print}) & New & Per-walk atom pools; load \(mk+D\cdot O(\lambda)\)\\
Algorithm, soundness, completeness, time, cap (\cref{sec:algorithm,sec:analysis,sec:asymptotics}) & New & Structure follows \cite[Sections~7--9]{NagCG}\\
Almost-everywhere transfer (\cref{sec:transfer}, \cref{app:transfer}) & Imported / adapted & \cite[Section~11 and Appendices~C--D]{NagCG}, from \cite{CLW20,CW19,BW24,MW20,GJS76}; wire measure in place of top weight, every class-dependent step re-derived in \cref{app:transfer}\\
\bottomrule
\end{tabularx}
\caption{Attribution of the components of the proof.}
\label{tab:attribution}
\end{table}


\subsection{Formal counterparts of the proof diagram}\label{app:lean-map}
\Cref{tab:lean-nodes} maps the nodes of \cref{fig:proof-dag} to
formal declarations. The adapter \path{Adapters.sources_to_nine} derives, from the six published results, the intermediate interfaces used by the proof (normalization, restriction, concentration, rectangular product, short PCP, refuter, amplification, three-query PCPP). Wire names are relative to
\path{NearQuarticWires}; source-adapter names are relative to
\path{NearCubicGates}. The theorem \path{Release.main_from_published}
proves the joint statement from \path{Close.PublishedInputs}, with
the common constants in \path{Statement.JointStatement4}.

\begingroup\footnotesize
\setlength\LTleft{0pt plus 1fill}\setlength\LTright{0pt plus 1fill}
\begin{longtable}{@{}>{\raggedright\arraybackslash}p{.27\textwidth}>{\raggedright\arraybackslash}p{\dimexpr.73\textwidth-2\tabcolsep\relax}@{}}
\caption{Formal counterparts of the proof diagram.}\label{tab:lean-nodes}\\
\toprule Diagram component & Exact formal declaration\\\midrule
\endfirsthead
\toprule Diagram component & Exact formal declaration\\\midrule
\endhead
Williams & \path{Close.PublishedInputs.williams_rectangularProductLeaf}\\[3pt]
BCGT13 & \path{Close.PublishedInputs.bcgt13_theorem11_1}\\[3pt]
Shoup/BCGT & \path{Close.PublishedInputs.shoup_findIrrPolyAlgorithm17}\\[3pt]
STV & \path{Close.PublishedInputs.stv_theorem29}\\[3pt]
CW proximity & \path{Close.PublishedInputs.cw19Proximity_lemma22Half}\\[3pt]
CW code & \path{Close.PublishedInputs.cw19Code_lemma26}\\[3pt]

Published-source adapter & \path{Adapters.sources_to_nine}\\[3pt]
Refuter & \path{Adapters.Assembly.field7}\\[3pt]
Short PCP & \path{Adapters.Assembly.Fixed.field6OrdinaryFixed}\\[3pt]
Amplification & \path{Adapters.Amplification.field8}\\[3pt]
Three-query PCPP/code & \path{Adapters.Assembly.field9Encoder}\\[3pt]
Product-star cells & \path{Cells.partition_average}; \path{Cells.mean_population}\\[3pt]
Bounded-matrix sampler & \path{Sampler.mean_traceMoment_le}; \path{Sampler.paper_fourth_root_budget}\\[3pt]
Cycle-sum identity & \path{Cycle.cycle_sum_identity}; \path{Cycle.conditional_cycle_sum_identity}\\[3pt]
Integer masks & \path{Mask.maskPoly_eq_mask}; \path{Mask.indicatorPolynomial_eq_mask}\\[3pt]
Localization/tomography & \path{Walks.chain_localization_failure_bound}; \path{Walks.signed_residue_recovery}\\[3pt]
Printing & \path{Cells.compiledTable_value}; \path{Cells.compiledTable_weightedDominance}\\[3pt]
Certified counter & \path{Program.counterExists4}\\[3pt]
Wire transfer & \path{Transfer.transfer4_official}; \path{Program.finalStatement4}\\[3pt]
Wires and gate outputs & \path{Statement.WiresAt}; \path{Statement.SymGatesAt}; \path{Statement.GatesAt}; \path{GateCorollary.jointStatement4_of_final}\\[3pt]
Joint conclusion & \path{Release.main_from_published}; \path{Statement.JointStatement4}\\
\bottomrule
\end{longtable}
\endgroup

